%%
% This is an Overleaf template for scientific articles and reports
% using the TUM Corporate Desing https://www.tum.de/cd
%
% For further details on how to use the template, take a look at our
% GitLab repository and browse through our test documents
% https://gitlab.lrz.de/latex4ei/tum-templates.
%
% The tumarticle class is based on the KOMA-Script class scrartcl.
% If you need further customization please consult the KOMA-Script guide
% https://ctan.org/pkg/koma-script.
% Additional class options are passed down to the base class.
%
% If you encounter any bugs or undesired behaviour, please raise an issue
% in our GitLab repository
% https://gitlab.lrz.de/latex4ei/tum-templates/issues
% and provide a description and minimal working example of your problem.
%%


\documentclass[
  english,        % define the document language (english, german)
  font=times,     % define main text font (helvet, times, palatino, libertine)
  twocolumn,      % use onecolumn or twocolumn layout
]{tumarticle}


% load additional packages
\usepackage{lipsum}


% article metadata
\title{Agentic AI for Code Generation}
\subtitle{Seminar Report}

\author[affil=1, email=thuaduc.nguyen@tum.de]{Student: "Thua Duc Nguyen"}

\author[affil=1, email=derui.zhu@tum.de]{Supervisor: "Derui Zhu"}


\affil[mark=1]{\theDepartmentName, \theUniversityName}



\date{April 1, 2025}


\begin{document}

\maketitle

\begin{abstract}
  % TODO: Write abstract
\end{abstract}

\section{Introduction}

\paragraph{Motivation.}
The rise of ``vibe coding'' and AI-assisted software development:
how natural-language-driven programming is reshaping the developer workflow.

\paragraph{Background and Related Surveys.}
Positioning relative to prior surveys on LLM-based code generation,
automated program repair, and AI agent architectures.
What this report adds: a pipeline-centric synthesis across four active research fronts.

\paragraph{Scope and Contribution.}
Define the scope (autonomous code-generation agents, 2023--2025 literature),
state the contribution (unified pipeline framing, cross-topic synthesis),
and outline the report structure.

\section{The Vibe Coding Pipeline}

Present the vibe coding pipeline as a unifying conceptual framework
and compare it to classical SE process models (waterfall, agile inner loops).

\subsection{Pipeline Stages}
Formal definition of each stage: Intent Capture \textrightarrow{} Planning \textrightarrow{}
Code Generation \textrightarrow{} Execution \textrightarrow{} Refinement.
Five stages forming a loop: execution feedback drives refinement, which may trigger
re-planning or re-generation. What inputs and outputs connect adjacent stages.

\subsection{Mapping to Survey Topics}
Explicit mapping: how each of the four subsequent sections corresponds to one or more
pipeline stages. \S3 covers the full pipeline as an integrated system;
\S4 zooms into planning; \S5 covers multi-agent approaches to generation and validation;
\S6 covers the execution-feedback-refinement loop (including fault-specific self-repair).
Cross-cutting concerns (evaluation, context management) are consolidated in the Discussion.

\section{End-to-End Autonomous Software Engineering Agents}
The big picture first: systems that orchestrate the entire pipeline---from a
natural-language issue to a tested pull request---showing how the stages
compose into a holistic system before we zoom into each stage.

\subsection{Architecture and Orchestration}
Common architectural choices: tool-augmented LLMs, ReAct-style loops, memory modules,
context management strategies, and how agents interface with codebases and dev tools.
Emphasis on the orchestration layer that ties pipeline stages together.

\subsection{Representative Systems}
Concrete systems (e.g., SWE-Agent, Devin, OpenHands, AutoCodeRover, Agentless)
with brief descriptions of their design and reported full-pipeline performance.

\subsection{Full-Pipeline Evaluation}
How end-to-end agents are evaluated: SWE-bench resolution rates, real-world issue benchmarks,
and the methodology behind current leaderboards. Critical analysis of benchmark
limitations is deferred to the Discussion (\S7.2).

\section{Planning and Search Strategies for Code Synthesis}
Zooming into the first critical stage: how agents decide \emph{what} to generate
before or during generation---decomposing problems, exploring solution spaces,
and selecting among candidates.

\subsection{Plan Representations and Decomposition}
Pseudocode plans, hierarchical task decomposition, chain-of-thought outlines,
dependency graphs, and specification sketches used to guide generation.

\subsection{Search and Exploration Strategies}
Tree/graph search over partial programs (MCTS, beam search, best-first),
sampling-and-ranking, evolutionary approaches, and backtracking mechanisms.

\subsection{Plan--Generate Interplay}
How planning interleaves with code generation: plan-then-generate vs.
interleaved planning, plan revision upon execution feedback, and plan grounding.

\section{Multi-Agent Frameworks for Code Generation}
An alternative architectural approach to the generation-to-validation span:
multiple specialized agents divide labor by role (architect, coder, tester, reviewer),
collectively covering planning, generation, and internal verification
to improve code quality beyond what a single agent achieves.

\subsection{Role Decomposition and Communication}
How tasks are split among agents, what roles exist, and how agents
exchange information (shared blackboard, message passing, structured dialogue).

\subsection{Framework Designs}
Survey of concrete frameworks (e.g., ChatDev, MetaGPT, AutoGen, CrewAI)
and their orchestration mechanisms.

\subsection{When Does Multi-Agency Help?}
Conditions under which multi-agent designs outperform monolithic agents;
coordination overhead, failure modes, and diminishing returns with agent count.

\section{Execution Feedback and Iterative Refinement}
Closing the loop: agents execute generated code, observe runtime outcomes
(test failures, compiler errors, exceptions), and revise their output.
Covers both general refinement and fault-specific self-repair.

\subsection{Feedback Signals and Environments}
Types of feedback: test results, compiler errors, runtime exceptions, static analysis,
human approval signals. Sandboxed execution environments and their design.

\subsection{Iterative Refinement Mechanisms}
How agents incorporate feedback: reflexion-style verbal self-critique, repair prompts,
diff-based editing, and multi-turn conversation with the environment.

\subsection{Fault Localization and Self-Repair}
When refinement becomes debugging: stack-trace analysis, LLM-driven root-cause reasoning,
edit-based patching, candidate ranking, and test-driven repair loops.

\subsection{Convergence and Stopping Criteria}
When to stop iterating: pass-rate thresholds, diminishing improvement detection,
budget limits, and the risk of over-fitting to weak test suites.

\section{Discussion}

\subsection{Consolidated Patterns and Trends}
Recurring design patterns across all five topics (observe-think-act loops, tool orchestration,
LLM-as-judge verification). Macro trends: from single-shot prompting to agentic scaffolds,
from monolithic agents to composable pipelines, from flat CoT to tree-structured search.

\subsection{The Evaluation Landscape}
Cross-cutting view of benchmarks (HumanEval, MBPP, SWE-bench, Defects4J, CrossCodeEval),
their coverage of pipeline stages, saturation concerns, and gaps in realism.

\subsection{The Context Window Bottleneck}
How finite context constrains every pipeline stage; retrieval-augmented generation,
hierarchical summarization, and repository-level chunking as mitigations.

\subsection{Open Challenges}
Unsolved problems: long-horizon reasoning, specification ambiguity,
cost efficiency, trustworthiness, and human-in-the-loop control design.

\subsection{Implications for Software Engineering Practice}
How these advances may reshape developer workflows, team structures,
and the economics of software production.

\section{Conclusion}
Summarize key findings across the five research fronts,
restate the pipeline framing as a unifying lens,
and point to the most promising future directions.

\end{document}
